\section{完整推理系统与实验方法}
\label{sec:system-method}

\subsection{完整双尺度推理}

系统执行 Ultralytics YOLOv3-tiny 的完整双尺度 DAG\cite{ultralyticsv95}。
主干卷积 m0、m2、m4、m6、m8、m10、m13、m14 和 m15 依次由 LASA 执行。
m8 的 26$\times$26$\times$256 输出同时作为高分辨率分支保存，并经普通
2$\times$2/s2 最大池化生成后续主干输入；m10 后执行右侧和底部补量化零值的
2$\times$2/s1 最大池化。m16 将主干分支压缩到 128 通道，经最近邻上采样后
与 m8 分支重量化到同一数值域，并按 HWC 顺序拼接为
26$\times$26$\times$384。m19 和两个 255 通道检测头随后仍由 LASA 执行。

\begin{figure*}[t]
\centering
\placeholderbox{50mm}{视觉占位：完整 13 卷积双尺度 DAG\\深色节点表示 PL 卷积，浅色节点表示 A53 池化、上采样、重量化拼接和双头后处理}
\caption{完整 YOLOv3-tiny 的软硬件划分。规则卷积全部映射到 LASA，低计算密度
且控制不规则的张量变换和检测后处理由 A53 完成。}
\label{fig:dag-placeholder}
\end{figure*}

四核 Cortex-A53 采用控制核与工作核分工。控制核维护 DAG 依赖、DMA 和 PL
上下文；工作核按独立空间或通道范围执行池化、上采样和分支重量化。检测头输出
首先生成候选框，再按类别执行非极大值抑制\cite{nms2017}。完整网络只在层边界
访问紧凑 HWC 张量，因此软件张量算子与 LASA 的输入输出布局一致，无需额外
转置。

\begin{table}[t]
\centering
\caption{A53 执行的非卷积节点及其数据作用。}
\label{tab:a53-ops}
\small
\begin{tabular}{lp{47mm}}
\toprule
节点 & 作用 \\
\midrule
普通池化 & 主干空间尺寸减半，同时保留 m8 高分辨率分支 \\
特殊池化 & 右/下补量化零值后执行 2$\times$2/s1 最大池化 \\
上采样 & 将 13$\times$13$\times$128 最近邻扩展到 26$\times$26 \\
重量化拼接 & 两分支统一数值域并形成 384 通道 HWC 张量 \\
双头后处理 & 候选生成、坐标恢复和类别独立 NMS \\
\bottomrule
\end{tabular}
\end{table}

选择这些节点在 A53 执行基于两个特征：其算术密度低于卷积，且控制或数据形状
由分支关系决定。四核按输出行或通道划分后，每个工作核写入互不重叠的目标区间，
控制核只在节点边界同步。分支重量化采用同一整数公式和饱和规则，保证 m19 看到
单一输入量化域。该软硬件边界保留了完整网络语义，也避免为少量不规则节点扩大
PL 控制范围。

输入可以由持久网络会话或 SD 冷启动程序送入同一内存接口。性能运行中，模型
参数常驻 DDR，输入已经位于内存，输出为完成双尺度后处理后的检测结果。数据
传输和文件系统时间单独统计，不进入 resident 延迟。该定义对应持续服务场景，
也使 PL、A53 张量算子和检测后处理能够独立归因。

\subsection{平台与数值配置}

实验平台为 AMD Xilinx KV260，器件为 XCK26，LASA 和 AXI DMA 工作在
200 MHz；设计使用 Vivado/Vitis 2022.2\cite{xilinxkv260}。阵列参数为
$R=18$、16 列双输出通道和 32 通道输出块。输入固定为 batch=1、
416$\times$416、RGB HWC。网络为 COCO80 双检测头 YOLOv3-tiny，评测数据为
COCO val2017 的 5000 张图像\cite{coco2014}。

卷积使用对称 INT8 权重和仿射 uint8 激活。主机整数参考直接模拟输入中心化、
INT8 乘法、INT32 累加和偏置、硬件舍入、饱和、激活查找表及所有 A53 张量
算子。该参考输出与参数包使用同一整数常量，因而节点比较检验的是部署语义，
而非浮点近似。模型精度分别报告相同 416 预处理下的 FP32 与 INT8 结果；
640 输入仅用于核对上游权重的公开基准。

\subsection{正确性、精度与性能口径}

正确性验证分为三个层次。第一层选择 128 张涵盖 80 类、不同长宽比和目标尺度
的图像，对 22 个整数节点逐字节比较，其中包含 13 个 PL 卷积及 A53 张量变换。
第二层对全部 5000 张图像比较两个原始检测头，并使用相同主机后处理计算 COCO
指标。第三层通过持续运行检查输入顺序、输出校验、DMA、PL状态和超时，验证
同一模型在长时间执行中的确定性。

精度采用 COCO bbox 评测，报告 AP50:95、AP50、AP75 以及小、中、大目标 AP。
部署比较的 FP32 和 INT8 路径消费相同的 416$\times$416 letterbox 字节与坐标
元数据。准确率模式使用较低候选阈值、类别独立抑制和最多 300 个输出，使网络
输出能够按照 COCO 标准评估。

性能实验独立运行 3 次，每次先预热 20 张，再计时 1000 张。resident 延迟从
量化输入已在 DDR 开始，到双尺度检测结果形成结束，包含全部 PL 卷积、A53
张量算子和检测后处理。计时区排除文件读写、网络传输和文本输出。结果报告均值、
P95、吞吐率、各层 PL 时间和 A53 分解，并通过硬件周期计数检查分项总和。

同板 CPU 对照在 KV260 的四核 Cortex-A53 上执行同一 13 卷积 INT8 DAG。
处理器工作频率为 \CPUFourCoreClockGHz{} GHz，卷积内核使用 NEON 向量指令并按
输出通道划分到四个核心；参数常驻 DDR，A53 执行的池化、上采样、重量化拼接
和双头后处理与 LASA 路径保持一致。该对照采用单张 416$\times$416 图像的
resident 计时，排除 JTAG 下载、存储和网络传输以及文本输出，并要求两个原始
检测头与整数参考逐字节一致。由此，同板比较只改变 13 个卷积的执行资源，不
改变网络、量化参数、输入或后处理语义。

计时器在任何文本输出之前停止。PL 层记录从四路 DMA 和加速器均进入可运行
状态开始，到输出 DMA 完成且加速器提交结束；A53 节点记录其并行区域和屏障；
resident 记录覆盖完整 DAG。对每张图还记录 13 层时间、A53 张量算子时间、
候选生成/NMS 时间和总时间。若分项之和与 resident 边界出现不可解释偏差，该
样本被标记为运行错误，不进入成功分布。

阵列利用率采用式\eqref{eq:utilization} 的模型运算量口径，避免将填充乘加计入
有效工作。尾轮和尾通道会消耗部分物理周期，却只按真实网络 MAC 计入分子；
因此该指标同时反映映射效率、数据供应和调度空泡。FPS 则以完整 resident 延迟
计算，包含 A53 非卷积节点和后处理。两项指标分别回答阵列是否连续工作和完整
推理能否实时执行。
